\begin{abstract}
Traditional dynamic model ensembling and model merging techniques typically rely on linear interpolation or flat blending of parameters or activations. However, we show that when merging low-rank projection operators (such as those representing specialized task-specific representation spaces), linear interpolation inherently violates the geometry of the projection manifold. This leads to projected coordinate collapse and severe representation distortion due to eigenvalue shrinkage. To resolve this, we propose \textbf{Continuous Riemannian-Geometric Homotopical Model Merging via Grassmannian Geodesic Blending (C-Lie-MM)}, a mathematically rigorous framework that treats task-specific projection bases as points on the Grassmannian manifold $\mathcal{G}(d, D)$. C-Lie-MM maps these bases onto the tangent space of a reference subspace using the Riemannian logarithm map, performs dynamic weighted ensembling of tangent matrices governed by a temperature-calibrated Gibbs routing policy, and projects the blended tangent vector back onto the manifold via the Riemannian exponential map. This ensures that the merged operator is always a mathematically correct, orthogonal projection matrix of rank $d$ on $\mathcal{G}(d, D)$. Empirically, evaluated inside a 14-layer Analytical Coordinate Sandbox simulation, C-Lie-MM completely prevents the representation collapse of uniform merging (which drops to 25.00\% accuracy under severe manifold overlap), achieving 70.30\% $\pm$ 4.01\% accuracy and outperforming prior flat ensembling baselines. Furthermore, because C-Lie-MM computes dynamic barycenters sample-wise in a single forward pass, it remains perfectly immune to heterogeneity collapse under mixed streaming workloads, establishing its systems-level viability.
\end{abstract}
